% Data dan metadata dokumen skripsi atau proposal UNESA.
% Seluruh data administratif dikelola di sini agar tidak perlu mengubah berkas .cls.

% Judul naskah dalam bahasa Indonesia dan bahasa Inggris.
% Judul bahasa Inggris digunakan untuk sampul dan halaman abstrak bahasa Inggris.
\Judul{Pengembangan \textit{Knowledge Graph-Based Retrieval Augmented Generation} dengan \textit{Hybrid Vector-Graph Retrieval} untuk Asisten Pencarian literatur Akademik}

\JudulEng{Development of \textit{Knowledge Graph-Based Retrieval Augmented Generation} with \textit{Hybrid Vector-Graph Retrieval} for Academic Literature Search Assistant}

% Identitas penulis.
\Nama{Rizky Yanuar Kristianto}
\NIM{22031554017}

% Program studi, tahun kelulusan atau sidang, dan fakultas.
% Fakultas membutuhkan dua argumen: nama lengkap dan singkatan resmi.
\ProgramStudi{Sarjana Sains Data}
\Programme{Bachelor of Data Science}
\Tahun{2026}

\Fakultas{Matematika dan Ilmu Pengetahuan Alam}{MIPA}
\Faculty{Mathematics and Natural Sciences}
\Universitas{Universitas Negeri Surabaya}
\Institution{State University of Surabaya}

% Pejabat program studi dan fakultas.
\KPS{Yuliani Puji Astuti, M.Si.}
\NIPkps{197807312006042001}

\Dekan{Prof. Dr. Wasis, M.Si.}
\NIPDekan{196712031993021001}

% Dosen pembimbing dan NIP.
% Format: {Pembimbing 1}{Pembimbing 2}. Kosongkan kurung kedua {} jika hanya satu pembimbing.
\Pembimbing{Ibnu Febry K., S.Kom., M.Sc., Ph.D.}
{Hasanuddin Al-Habib, M.Si.}

\NIPPembimbing{198802182015041001}
{199308092022031010}

% Dosen penguji dan NIP.
% Format: {Ketua Penguji}{Anggota 1}{Anggota 2}. Kosongkan kurung ketiga {} jika hanya dua penguji.
\Penguji{Dr. Atik Wintarti, M.Kom.}
{Yuni Rosita Dewi, M.Si.}
{}

\NIPPenguji{196610121991032002}
{199506282024062001}
{}

% Pembimbing ahli atau mitra luar instansi (isi tanda strip jika tidak ada).
\Ahli{-}
\NIPahli{-}

% Tanggal persetujuan, ujian sidang, dan pernyataan orisinalitas.
\TanggalUji{3 juli 2026}
\TanggalDisetujui{17 juni 2026}
\TanggalSidang{3 Juli 2026}
\TanggalOrisinalitas{17 Juli 2026}

% Deklarasi penggunaan kecerdasan buatan pada surat pernyataan AI.
% Gunakan \checkedbox untuk mencentang atau \emptybox untuk membiarkan kotak kosong.
% Jika mencentang opsi lainnya, isi rincian penggunaannya pada \AILainnyaTeks.
\renewcommand{\AIEksplorasiIde}{\checkedbox}
\renewcommand{\AIKerangkaTulisan}{\checkedbox}
\renewcommand{\AITataBahasa}{\emptybox}
\renewcommand{\AIAnalisisData}{\checkedbox}
\renewcommand{\AIIlustrasi}{\emptybox}
\renewcommand{\AILainnya}{\emptybox}
\renewcommand{\AILainnyaTeks}{}
